\gddchapter{Testing sequence}
The prototype can be played with both voice and keyboard and the navigation modes can be swapped with a button press. 


\section{Novel version B - Voice navigation}

\subsection{Observation log}
\begin{table}[H]
\centering
\begin{tabular}{|p{5cm}|p{10cm}|}
\hline
\textbf{Field} & \textbf{Response} \\ \hline
\textbf{Physical comfort} How is the participant feeling & The player appears to be comfortable.  \\[1.5cm] \hline
\textbf{Natural language observation} (is the player using long sentences or simple commands? What type of language do they use? Are they speaking first or 3rd person?) & The participant is using direct command language "Go to X" "Do Y". \\[1.5cm] \hline
\textbf{Linguistic guessing} (is the particpant trying to guess the possible commands by speaking to the game?) & The player is making guesses and asking the researcher about the system boundaries \\[1.5cm] \hline
\textbf{Processing time user reaction} (Does the wait time break the suspension of disbelief) & There delay is very noticable to the player \\[1.5cm] \hline
\textbf{Voice processing pipeline edge cases} (System edge cases eg. two vs to)& There were no edge cases detected in his experience. \\[1.5cm] \hline
\end{tabular}
\end{table}


\section{Baseline A - Keyboard navigation}
After playing the audio version of the game, the player is now experiencing the baseline version with keyboard usage.
This allows the player to anchor in the typical experience and serves as a comparative point in the A/B testing.

\subsection{Observation log}
\begin{table}[h]
\centering
\begin{tabular}{|p{5cm}|p{10cm}|}
\hline
\textbf{Field} & \textbf{Response} \\ \hline
\textbf{Physical comfort} How is the participant feeling & The participant seems to be more comfortable in the keyboard version of the experience\\[1.5cm] \hline
\textbf{Immersion markers} (flow \& frustration) & Kuba seems to be more immersed in the keyboard version of the game.\\[1.5cm] \hline
\end{tabular}
\end{table}


\subsection{Sticky moments}
The moments that stood out during testing marked with timestamps.

\begin{table}[h]
\centering
\begin{tabular}{|p{3cm}|p{10cm}|}
\hline
\textbf{Timestamp} & \textbf{Log} \\ \hline
02:48 & Start of the playtest \\ \hline 
03:20 & Reading too much \\ \hline 
04:20 & What's the floor made of? \\ \hline 
06:00 & Scan the room \\ \hline 
08:30 & Look around for food \\ \hline 
\end{tabular}
\end{table}

\section{Alignment check}
How much was the researcher participating in the gameplay experience? 

The researcher did his best not to participate actively during the testing sequence. 



